\section{Model and inferential targets}
\label{sec:background}

\subsection{A branching experiment with a known observation scale}
\label{sec:bg-model}

Consider one generation of transmission with $n\ge1$ primary infections. Conditional on the reproduction number $R_t>0$, primary case $i$ produces independent negative binomial offspring \citep{lloydsmith2005superspreading},
\begin{equation}
  X_{i,t}\mid R_t\sim\NB\!\left(k,\frac{k}{k+R_t}\right),\qquad
  \E(X_{i,t}\mid R_t)=R_t,\quad
  \Var(X_{i,t}\mid R_t)=R_t+\frac{R_t^2}{k}.
  \label{eq:offspring}
\end{equation}
Here $k>0$ controls individual heterogeneity. Each secondary infection is independently reported with probability $\rho\in(0,1]$. Summation and binomial thinning give
\begin{equation}
\begin{aligned}
  C_t\mid n,R_t&\sim\NB\!\left(nk,\frac{k}{k+\rho R_t}\right),\\
  \E(C_t\mid n,R_t)&=\rho nR_t,\qquad
  \Var(C_t\mid n,R_t)=\rho nR_t+\frac{\rho^2nR_t^2}{k}.
\end{aligned}
  \label{eq:observed}
\end{equation}
Define $m=\rho n$ and $\Rhat_t=C_t/m$. The design scale $m$ is the expected reported count at $R_t=1$, not the realised count $C_t$.

The baseline assumptions are fixed, known $\rho$ and independent reporting (H1); conditionally independent offspring as in Eq.~(\ref{eq:offspring}) (H2); fixed $k$, known for the likelihood calculation (H3); and a fixed or exogenous $n$ independent of transmission (H4). Results are conditional on the selected design, with the same $m$ in each period. Write $\E(R_t)=\Rbar>0$ and $\Var(R_t)=v_R<\infty$. Across repeated periods we additionally assume that counts are independent given the entire transmission trajectory, with conditional law~(\ref{eq:observed}). These are repeated branching experiments; they do not assert that naturally evolving incidence is exogenous.

This observation scale matters. Stability of an unknown reporting fraction does not establish unbiasedness: dividing by an assumed $\rho_0$ yields conditional mean $(\rho/\rho_0)R_t$. More cases do not remove the resulting multiplicative bias. A ratio constructed from noisy observed incidence also has a random denominator and need not share the properties of $C_t/(\rho n)$. The baseline therefore specifies an informative reference experiment, not a direct model for every surveillance statistic.

The conditional counting variance includes both offspring randomness and independent omission. Conditioning first on the total secondary infections separates their contributions as $(R_t+R_t^2/k)/n$ and $(1-\rho)R_t/(\rho n)$. Their sum equals the variance above. In particular, the term linear in $R_t$ is not exclusively reporting noise: it persists with complete reporting. This distinction prevents a purely algebraic partition from being mistaken for a unique biological attribution.

\subsection{Current values, realised means and process means}
\label{sec:bg-variance}

For the current target, $\E(\Rhat_t\mid R_t)=R_t$ and the conditional variance is $[R_t+\rho R_t^2/k]/m$. It attains the conditional Cram\'er--Rao bound in the one-parameter count model (Appendix~\ref{app:proofs}). The statement concerns estimators unbiased for every $R_t$ under that likelihood. It is not a universal minimax bound, and it does not cover biased, hierarchical or prior-informed procedures.

\begin{proposition}[One period, two targets]
\label{prop:average}
Under H1--H4, let $A=\Rbar+\rho(\Rbar^2+v_R)/k$. Then
\begin{equation}
\begin{aligned}
  \E[(\Rhat_t-R_t)^2]&=\frac{A}{m},\\
  V(m):=\E[(\Rhat_t-\Rbar)^2]&=v_R+\frac{A}{m}.
\end{aligned}
  \label{eq:var-average}
\end{equation}
Among estimators based on this period's count that are conditionally unbiased for $R_t$, $\Rhat_t$ minimises both risks.
\end{proposition}

The first risk averages the conditional estimation error over the stated distribution of $R_t$. It is not the conditional risk at every fixed value, nor a posterior credible interval for a particular period. The second evaluates the same observation against a different target. No assumption of independent transmission periods is needed for this single-period comparison, because only the marginal moments enter. Independence becomes relevant when risks are aggregated across periods.

The distinction follows from
\begin{equation}
  \Rhat_t-\Rbar=(\Rhat_t-R_t)+(R_t-\Rbar).
  \label{eq:error-split}
\end{equation}
Conditional unbiasedness makes the two terms uncorrelated. Between-period transmission variation is part of the current signal, but contributes sampling error for the process mean. The latter is a mean under a specified stationary law, not automatically a seasonal average in a nonstationary epidemic. The two contributions have the familiar within-unit and between-unit structure of clustered sampling \citep{cochran1977sampling, engen1998demographic}.

For $v_R>0$, define the balance scale $\mx=A/v_R$. Increasing $m$ through $n$ at fixed $\rho$, $k$ and transmission law leaves $A$ unchanged: the two terms in $V(m)$ are equal at $m=\mx$. This scale describes the selected estimator and target. It does not mark a loss of all information, and changing $\rho$ changes $A$ as well as $m$. Coverage and budget calculations are given only as symbolic extensions in Appendix~\ref{app:extensions}.

\begin{corollary}[A fixed set of observed periods]
\label{cor:targets}
Under H1--H4 and conditional independence of counts given $R_{1:T}$, define $\Rbar_T=T^{-1}\sum_{t=1}^T R_t$ and $\widehat{\Rbar}_T=T^{-1}\sum_{t=1}^T\Rhat_t$. For integer $T\ge1$, with common marginal moments of $R_t$,
\begin{equation}
\begin{aligned}
  \E[(\widehat{\Rbar}_T-\Rbar_T)^2]&=\frac{A}{Tm},\\
  \E[(\widehat{\Rbar}_T-\Rbar)^2]&=B_T+\frac{A}{Tm},
  \qquad B_T=\Var(\Rbar_T).
\end{aligned}
  \label{eq:target-means}
\end{equation}
For independent, identically distributed $R_t$, $B_T=v_R/T$. For a stationary process with $\Cov(R_t,R_{t+h})=v_R\phi^{|h|}$, $|\phi|<1$, the exact finite-$T$ expression is $B_T=v_RS(\phi,T)/T^2$, where
\begin{equation}
  S(\phi,T)=T+2\sum_{h=1}^{T-1}(T-h)\phi^h
  =\frac{T(1+\phi)}{1-\phi}-\frac{2\phi(1-\phi^T)}{(1-\phi)^2}.
  \label{eq:ar-sum}
\end{equation}
\end{corollary}

The realised-mean identity does not require independent transmission levels: only the counting errors must be conditionally independent, with the stated common moments. Increasing cases for fixed $T$ removes its counting error, whereas process-mean risk approaches $B_T$. Additional periods address that term; dependence determines how much they contribute. Forecasting a future value is a further target, treated in Appendix~\ref{app:extensions}.

The stationary covariance formula uses integer periods and the complete finite sum, rather than a long-series approximation. Positive persistence usually reduces the gain from adjacent periods; negative autocorrelation can increase it. Neither a serially correlated process nor a season with changing mean should be replaced automatically by independent draws. For nonstationary transmission, the realised mean remains well defined, but interpretation of a common process mean requires an additional model.

Spatial averaging has an analogous condition. For $J$ units with equal transmission variance $v_R$ and mean pairwise correlation $\bar c$, the variance of their equally weighted transmission mean is
\begin{equation}
  \Var\!\left(J^{-1}\sum_{j=1}^J R_{j,t}\right)
  =v_R\left[\bar c+\frac{1-\bar c}{J}\right].
  \label{eq:pooling}
\end{equation}
This covariance identity is not a decomposition of an empirical regression intercept. Unequal variances, unequal weights or a changing correlation structure require the corresponding covariance matrix.

\subsection{Unobserved reporting volatility and identification}
\label{sec:bg-reporting}

Replace H1 by a common period-specific probability $p_t\in[0,1]$, independent of $R_t$ and the design, with known mean $\rho>0$ and $\Var(p_t)=\rho^2\tau^2$ (H5). Reporting remains independent across infections conditional on $p_t$. Necessarily $\tau^2\le(1-\rho)/\rho$. The estimator still divides by $m=\rho n$; it does not observe $p_t$.

\begin{proposition}[Reporting volatility and moment nonidentification]
\label{prop:report-noise}
Under H2--H5, put $Q=\Rbar^2+v_R$. Then
\begin{equation}
\begin{aligned}
  \E(\Rhat_t)&=\Rbar,\qquad
  V_H(m):=\Var(\Rhat_t)=v_R+v_H+\frac{A_H}{m},\\
  v_H&=\tau^2Q,\qquad
  A_H=\Rbar+\frac{\rho(1+\tau^2)Q}{k}.
\end{aligned}
  \label{eq:vH-AH}
\end{equation}
For known $m$, $\rho$ and $k$, the first two count moments depend on $v_R$ and $\tau^2$ only through $U=(1+\tau^2)(\Rbar^2+v_R)$. They therefore do not generally identify transmission variance and reporting volatility separately.
\end{proposition}

Conditional on both $R_t$ and $p_t$, the estimator's mean is $(p_t/\rho)R_t$. Unbiasedness for $R_t$ is restored only after averaging over independent reporting with the correct mean. This explains why more infections within a period cannot average away a shared reporting shock. If that period's reporting probability were measured and positive, a correctly adjusted estimator would pose a different problem; uncertainty in measuring that probability would itself require propagation.

This is a moment nonidentification result within the specified family. It does not say that external reporting measurements, temporal information under additional assumptions, or a more restrictive full distribution could never identify the mechanisms. Boundary cases can also impose further restrictions. The independence of reporting and transmission is substantive; if illness pressure changes reporting practice, the displayed decomposition need not hold.

\begin{corollary}[Current and realised-mean risks with volatile reporting]
\label{cor:report-current}
Under H2--H5,
\begin{equation}
  \E[(\Rhat_t-R_t)^2]=v_H+\frac{A_H}{m}.
  \label{eq:vH-current}
\end{equation}
If, in addition, the $p_t$ are independent and identically distributed, independent of the entire transmission trajectory, and counts are independent conditional on both trajectories, then
\begin{equation}
  \E[(\widehat{\Rbar}_T-\Rbar_T)^2]=\frac{v_H+A_H/m}{T},\qquad
  \E[(\widehat{\Rbar}_T-\Rbar)^2]=B_T+\frac{v_H+A_H/m}{T}.
  \label{eq:report-means}
\end{equation}
\end{corollary}

Thus, absence of a transmission-variation term for the realised target does not imply absence of reporting error. At fixed $T$, unobserved reporting volatility remains as $m$ grows. Serially dependent reporting adds covariance terms, and a stable unknown mean reporting fraction can add bias. Table~\ref{tab:targets} collects the identities and their conditions.

\begin{table}[t]
\centering
\small
\setlength{\tabcolsep}{4pt}
\caption{Inferential targets and exact mean squared errors. The fixed-reporting column assumes H1--H4. The volatile-reporting column assumes H2--H5 and, for averages, independent reporting probabilities independent of the entire transmission path. Counting errors are conditionally independent across periods. $B_T=\Var(\Rbar_T)$ equals $v_R/T$ for independent, identically distributed transmission and $v_RS(\phi,T)/T^2$ under the stated stationary autocovariance.}
\label{tab:targets}
\begin{tabularx}{\linewidth}{>{\raggedright\arraybackslash}X>{\centering\arraybackslash}p{0.16\linewidth}>{\centering\arraybackslash}p{0.24\linewidth}>{\centering\arraybackslash}p{0.27\linewidth}}
\toprule
Target & Estimator & Fixed, known reporting & Unobserved volatile reporting \\
\midrule
Current $R_t$ & $\Rhat_t$ & $A/m$ & $v_H+A_H/m$ \\
Realised mean $\Rbar_T$ & $\widehat{\Rbar}_T$ & $A/(Tm)$ & $(v_H+A_H/m)/T$ \\
Process mean $\Rbar$ & $\widehat{\Rbar}_T$ & $B_T+A/(Tm)$ & $B_T+(v_H+A_H/m)/T$ \\
\bottomrule
\end{tabularx}
\end{table}

\subsection{Endogenous count scales and apparent statistics}
\label{sec:bg-scope}

If $m=\rho n$ is random and associated with transmission, even retaining the conditional count law and unbiasedness given $(R_t,m)$ gives
\begin{equation}
\begin{aligned}
 \E[(\Rhat_t-\Rbar)^2\mid m]
 &=\frac{\E(R_t\mid m)+\rho\E(R_t^2\mid m)/k}{m}\\
 &\quad+\Var(R_t\mid m)+[\E(R_t\mid m)-\Rbar]^2.
\end{aligned}
\label{eq:endogenous-risk}
\end{equation}
These conditional moments can change with scale, so a constant intercept plus a constant inverse-count coefficient does not follow. Conditioning on an observed outcome count is still more restrictive: it need not preserve the assumed conditional observation law. A renewal representation alone does not resolve either problem.

The empirical response is consequently defined as squared relative deviation of an apparent statistic from its specified centre (Section~\ref{sec:meth-stats}), not error against known $R_t$. Its affine intercept and slope need not equal $v_R/\Rbar^2$ and $A/\Rbar^2$. Delays, aggregation, denominator noise and changing conditional means all affect that mapping. Calibration specifies a model-dependent connection (Section~\ref{sec:meth-inversion}); it does not turn the empirical intercept into an exact additive decomposition of biological sources.
